\HMTNarrativeHeading{Ce que couple la constante}

La constante holographique de couplage arithmético-géométrique est introduite
là où son registre est produit. Sa représentation entière est notée
\(K\) ; la lecture périodique, \(\kappa_{\rm ag}\). L’origine de
phase, l’orientation et l’ordre des composantes font partie de la
définition. Le mot « constante » nomme l’objet déterminé par cette
construction et par ses lecteurs, qui seront exposés au moyen d’opérations
explicites.

L'extraction réunit les canaux orientés et une donnée de somme. L'inversion de Hadamard recompose leurs composantes dans l'ordre naturel. Chaque coefficient de cette inversion provient de l'opérateur déclaré et de son image entière. Ainsi, le registre est déterminé avant son utilisation dans la clôture de la constante de structure fine. La lecture positionnelle transforme ensuite les douze composantes en une expression rationnelle périodique. La base et la longueur étant fixées, cette conversion admet un inverse exact qui restitue tous les blocs.

Le même ordonnancement a une fonction géométrique. Ses lecteurs distinguent
un quadruplet marqué et une configuration d’incidence ; la signature associée
à alpha sélectionne une hexade dans le plan correspondant. La
relation avec les régions numériques utilise les applications
d’alignement et de relèvement qui seront démontrées dans la partie exceptionnelle.
La lecture périodique et la configuration incidentielle reçoivent ainsi une
source commune, avec des opérations de types différents.

Le développement opératoire ultérieur ajoute une troisième fonction. Les
décalages cycliques du registre forment une base de son porteur. Les
réponses sur cette base déterminent un opérateur, et la forme de Gram
fournit une borne de stabilité pour le retrouver. Lorsque la
symétrie cyclique spécifiée existe, une réponse vectorielle complète permet de
reconstruire les autres. L’archive isométrique conserve cette capacité
d’identification.

Tel est le contenu du couplage qui structure l'exposé : codage arithmétique réversible, incidence marquée et identification opératorielle. On construit d'abord le registre ; on prouve ensuite ce que chaque lecteur en obtient. La théorie de récupération sera développée lorsque ses espaces et ses transports seront disponibles. Cette position explique conjointement la centralité de la constante et la hiérarchie des démonstrations qui utilisent ses différentes réalisations.

